\section{Introduction}
\label{sec:introduction}

A research plan contains more than a topic and a collection of well-formed sentences. Its central problem, objectives, literature topics, planned methods, and expected outputs need intelligible relationships if the plan is to be examined and revised. A practical difficulty motivating the Research Map is that these components can be drafted separately without making their connections clear. Sergio Fran\c{c}a's teaching map and later methodological specification address that difficulty by relating five planning stages to a companion objective--method matrix and by calling for human validation of proposed content \cite{franca2026map,franca2026detail}. This motivation identifies a design problem; it does not establish how often researchers encounter it or whether a structured map reduces delays or improves their work.

AI-supported research tools already address adjacent tasks. Published systems assist with staged proposal development, formulation of research questions, generation and critique of research ideas, and inspection of revision histories \cite{guerard2026researchllm,wang2026fivestage,liu2024coquest,baek2025researchagent,freire2006vistrails}. Publicly documented planning applications also describe persistent records, review checkpoints, or revision assistance \cite{projectready2026guide,researchly2026site}. These antecedents make generic claims about an agent, a sequence of prompts, or human oversight insufficient to distinguish the present work. The question for this article is more specific: \emph{How are the Research Map's methodological relationships represented and supported in a human-supervised AI application, and where does that representation remain incomplete?}

We address the question as a method-and-system account. First, we describe the Research Map from its two author-supplied methodological documents, separating the five numbered stages from the accompanying matrix and from later proposals for software support. Second, we examine the inspected implementation: editable AI proposals, confirmed planning units, links, revisions, checks, conditional advisor review, and export. The unit of analysis is the correspondence between methodological relations and software structures, including mismatches between them. Our contribution is a documented, inspectable account of the correspondence between a particular planning method and its current application, together with the limits of that correspondence. We do not claim that its individual functions are unprecedented or that the software validates a research design scientifically.

The evidence combines the primary method documents, a focused narrative comparison with scholarly work and publicly described applications, inspection of relevant source code and local technical checks, and one constructed walkthrough. The walkthrough uses a fictional library project to show how an upstream change could prompt reconsideration of objectives, topics, and the matrix. It is an explanatory example, not an application session or a study with participants. No user outcomes, methodological ratings by the authors, time savings, or comparative effectiveness are reported. This scope permits a critical account of what the method specifies and what the inspected system supports, while leaving questions of use and impact for separate empirical work.

Section~\ref{sec:state-of-the-art} positions the work among research-planning agents, provenance systems, and public applications. Section~\ref{sec:method} presents the Research Map method; Section~\ref{sec:system} examines its operationalization in the application and agent. Section~\ref{sec:walkthrough} gives the constructed example, followed by the contribution boundary and limitations in Section~\ref{sec:discussion} and the conclusion in Section~\ref{sec:conclusion}.
